\vspace{0.5em}
\noindent\textbf{References.} Each group matches the pitch section it supports.
\begingroup\sloppy

\vspace{0.3em}
\noindent Section 2, Goal. Existing programmable iris builds.
\begin{itemize}
    \item[{[G1]}]\phantomsection\label{ref:G1} PetaPixel, ``DIY digital iris animates and customizes lens bokeh'', 11 February 2026. Build report on an LCD iris fitted to a camera lens. \\
    \url{https://petapixel.com/2026/02/11/diy-digital-iris-animates-and-customizes-lens-bokeh/}
    \item[{[G2]}]\phantomsection\label{ref:G2} Adafruit Blog, ``Turning a vintage Nokia display into a programmable camera iris'', 10 February 2026. The same build from the hardware side, including how the builder stripped the panel. \\
    \url{https://blog.adafruit.com/2026/02/10/turning-a-vintage-nokia-display-into-a-programmable-camera-iris/}
\end{itemize}

\noindent Section 2, Goal. Programmable aperture research.
\begin{itemize}
    \item[{[G3]}]\phantomsection\label{ref:G3} C.-K. Liang et al., ``Programmable aperture photography: multiplexed light field acquisition'', ACM SIGGRAPH 2008. Source for the sub-aperture scan and the S-matrix multiplexing behind post-capture refocus. ResearchGate record. \\
    \url{https://www.researchgate.net/publication/220184327_Programmable_aperture_photography_Multiplexed_light_field_acquisition}
    \item[{[G4]}]\phantomsection\label{ref:G4} O. Bimber, H. Qureshi, A. Grundh\"ofer, M. Grosse and D. Danch, ``Adaptive coded aperture photography'', in Advances in Visual Computing, Lecture Notes in Computer Science, Springer, 2011. The paper changes the aperture pattern per frame to suit the scene, which forms the basis of the dynamic aperture mode. \\
    \url{https://link.springer.com/chapter/10.1007/978-3-642-24028-7_6}
    \item[{[G5]}]\phantomsection\label{ref:G5} Columbia CAVE laboratory, project page ``Programmable Imaging: Controllable Apertures''. Overview of controllable aperture optics and the results they produce. \\
    \url{https://cave.cs.columbia.edu/projects/categories/project?cid=Computational+Imaging&pid=Programmable+Imaging+Controllable+Apertures}
    \item[{[G6]}]\phantomsection\label{ref:G6} C. Zhou, S. Lin and S. K. Nayar, ``Coded aperture pairs for depth from defocus'', IEEE International Conference on Computer Vision (ICCV), 2009. Source of the two coded patterns and the joint deconvolution behind the priority goal.
\end{itemize}

\endgroup
